\section{Organisation du rendu}

\begin{UPSTIinfor}{Livrables attendus}
Votre rapport d'équipe (3 personnes) devra rassembler l'ensemble des
productions demandées dans ce sujet :
\begin{enumerate}
    \item Le schéma synoptique de l'installation ;
    \item Le schéma de l'architecture réseau, avec les adresses IP des
    différents équipements ;
    \item L'étude des modes de fonctionnement (GMMA) de la cellule et de
    l'automate, avec conditions d'entrée/sortie et actions pour chaque
    mode ;
    \item La cartographie mémoire de l'automate M340 pour les échanges avec
    le robot, l'IHM et le convoyeur ;
    \item Les sécurités mises en œuvre ;
    \item Les trajectoires du robot envisagées, avec leurs types de
    mouvements et les vitesses associées.
\end{enumerate}
\end{UPSTIinfor}

\begin{UPSTIidee}{Conseil}
Répartissez le travail au sein de votre équipe en fonction du planning
indicatif donné en introduction, mais prenez le temps, en fin de séance, de
relire l'ensemble du rapport à trois : la cohérence entre les livrables
(adressage réseau utilisé dans la cartographie mémoire, modes cohérents avec
les sécurités, etc.) sera un critère important d'évaluation.
\end{UPSTIidee}
